% sections/related.tex
% Included from main.tex via \input{sections/related}

\section{Related Work}
\label{sec:related}

Simulation platforms relevant to autonomous systems span autonomous driving,
aerial robotics, joint co-simulation, and embodied AI.
From the perspective of air-ground embodied intelligence, the central question
is not whether a platform supports driving or flight in isolation, but whether
aerial and ground agents can be jointly simulated within a unified, physically
coherent, and practically usable environment.
As illustrated in Fig.~\ref{fig:platform_positioning} and
Table~\ref{tab:platform_comparison}, existing open-source platforms largely
remain separated by domain focus, and none simultaneously provides realistic
urban traffic, socially-aware pedestrians, physics-based multirotor flight,
preserved native APIs, and single-process execution in one shared simulation
world.

\subsection{Autonomous Driving Simulators}
\label{sec:related:driving}

Autonomous driving simulators provide strong support for realistic urban
scenes, traffic agents, and ground-vehicle perception.
CARLA~\cite{carla}, built on Unreal Engine~\cite{ue4}, has become the de facto
open-source platform for urban driving research due to its photorealistic
environments, rich actor library, and mature Python API.
LGSVL~\cite{lgsvl} offers full-stack integration with Autoware and Apollo on
the Unity engine.
SUMO~\cite{sumo} provides lightweight microscopic traffic flow modeling.
MetaDrive~\cite{metadrive} enables procedural environment generation for
generalizable RL, and VISTA~\cite{vista} supports data-driven sensor-view
synthesis for autonomous vehicles.
These platforms collectively cover a broad range of ground-autonomy research
needs, but none natively supports physics-based UAV flight, leaving air-ground
cooperative workloads outside their scope.

\subsection{Aerial Vehicle Simulators}
\label{sec:related:uav}

Aerial simulators provide the complementary capability: accurate multirotor
dynamics, onboard aerial sensing, and UAV-oriented control interfaces.
AirSim~\cite{airsim} remains one of the most widely adopted open-source UAV
simulators, offering physics-accurate multirotor flight and a comprehensive
sensor suite on Unreal Engine, though its upstream development has since been
archived.
Flightmare~\cite{flightmare} combines Unity-based photorealistic rendering
with highly parallel dynamics for fast RL training.
FlightGoggles~\cite{flightgoggles} provides photogrammetry-based environments
for perception-driven aerial robotics.
Gazebo~\cite{gazebo}, together with MAV-specific packages such as
RotorS~\cite{rotorsim}, offers a mature ROS-integrated simulation stack for
multi-rotor control and state estimation.
OmniDrones~\cite{omnidrones} and
gym-pybullet-drones~\cite{pybulletdrones} target scalable, GPU-accelerated or
lightweight RL-oriented multi-agent UAV training.
While these systems are well suited to aerial autonomy in isolation, they
generally lack realistic urban traffic populations, pedestrian interactions, and
richly populated ground environments, limiting their use as infrastructure for
air-ground cooperation or cross-domain data collection.

\subsection{Joint and Co-Simulation Platforms}
\label{sec:related:joint}

The most relevant prior efforts attempt to combine aerial and ground simulation
through co-simulation.
TranSimHub~\cite{transimhub} connects CARLA with SUMO and aerial agents via a
multi-process architecture supporting synchronized multi-view rendering.
Other representative approaches include ROS-based pairings of AirSim with
Gazebo~\cite{airsim,gazebo,ros2}.
These systems demonstrate that heterogeneous simulation backends can be
functionally connected, but their integration typically depends on bridges, RPC
layers, or message-passing middleware across independent processes.
As summarized in Table~\ref{tab:related_arch}, such designs do not preserve a
single rendering pipeline, do not provide strict shared-tick execution, and
often require adapting existing code to new interfaces.
By contrast, CARLA-Air integrates both simulation backends within a single
Unreal Engine process, preserving both native APIs while maintaining a shared
world state, shared renderer, and synchronized sensing---a system-level
distinction detailed in Section~\ref{sec:architecture}.

\subsection{Embodied AI and Robot Learning Platforms}
\label{sec:related:embodied}

Embodied AI platforms prioritize a different design objective: scalable policy
training rather than realistic urban air-ground infrastructure.
Isaac Lab~\cite{isaaclab} and Isaac Gym~\cite{isaacgym} emphasize massively
parallel GPU-accelerated reinforcement learning for locomotion and
manipulation.
Habitat~\cite{habitat} and SAPIEN~\cite{sapien} target indoor navigation and
articulated object interaction, while RoboSuite~\cite{robosuite} focuses on
tabletop manipulation benchmarks.
These platforms are valuable for embodied intelligence research, but they do
not provide the urban traffic realism, socially-aware pedestrian populations,
or integrated aerial-ground simulation required by low-altitude cooperative
robotics.
In this sense, they address complementary research needs and are not direct
substitutes for CARLA-Air.

\subsection{Summary}
\label{sec:related:summary}

Fig.~\ref{fig:platform_positioning} positions CARLA-Air and representative
platforms along two principal design axes: simulation fidelity and agent domain
breadth.
Driving simulators provide realistic urban ground environments without aerial
dynamics; UAV simulators provide flight realism without populated ground
worlds; joint simulators generally rely on multi-process bridging that
sacrifices interface compatibility or synchronization fidelity; and embodied AI
platforms focus on scalable learning rather than air-ground infrastructure.
CARLA-Air is designed to sit at the intersection of these domains, combining
realistic urban traffic, socially-aware pedestrians, physics-based multirotor
flight, preserved native APIs, and single-process execution within one unified
simulation environment.
Table~\ref{tab:platform_comparison} provides a detailed feature-level comparison
across all platforms discussed above.

\input{figures/fig_platform_positioning}
\input{figures/tab_platform_comparison}
\input{figures/tab_related_arch}